\section{Discussion}

\subsection{A diagnostic rather than a mechanism test}

The main outcome of this study is a diagnostic decomposition of finite-sample
multifractal evidence. A high scaling score means that scaling is stable over
the tested scales. It does not imply nonlinear curvature. A nonzero
\(\lambda^2_{\mathrm{proj}}\) means that the empirical spectrum has an
MRW-family projection coordinate. It does not prove that the sample was generated
by an MRW. A high MRW diagnostic score should be interpreted only together
with the monofractal residual, MRW residual, curvature score, boundary score,
and instability warnings.

This caution is not a weakness of the framework; it is the point of the
framework. Finite empirical spectra often invite overinterpretation. CMIN-SR
is designed to make such overinterpretation harder by forcing the analyst to
compare competing explanations.

\subsection{Relation to classical multifractal estimators}

Classical estimators such as structure functions, MFDFA, and wavelet-based
methods remain central to multifractal analysis
\cite{muzy1991wavelets,muzy1993multifractal,kantelhardt2002mfdfa,jaffard2007wavelet,wendt2007bootstrap}.
CMIN-SR is
not proposed as a replacement for these methods. It is better understood as a
spectral diagnostic layer that can sit on top of empirical spectra and ask how
the estimated geometry relates to monofractal and MRW-family projections.

The supplemental baseline results are consistent with this interpretation.
MFDFA can reduce error relative to a simple structure-function baseline in
some settings, but high-\(\lambda^2\) MRW curvature remains difficult in the
finite ensemble baseline. The added Haar wavelet-leader and WTMM-style
controls provide a lightweight wavelet comparison and do not remove the
short-window recovery limitation. A full production WTMM or wavelet-leader
package would still be valuable as future benchmarking, but the present
classical controls already support the conservative finite-sample conclusion.

\subsection{Implications for complex systems}

The framework is especially relevant for complex systems where multiple
mechanisms can produce similar spectral signatures. In finance, for example,
heavy tails, volatility clustering, regime changes, and temporal dependence
can all affect scaling statistics
\cite{cont2001stylized,bouchaud2000apparent,zhou2009components,barunik2012source}.
The Fama-French sanity check shows small positive original-vs-shuffled gaps,
which is consistent with some temporal multiscale dependence. But this is not
MRW mechanism validation. The right conclusion is narrower: the diagnostic can
be applied to real return series and shows original-vs-shuffled differences
consistent with temporal structure.

\subsection{Limitations}

Several limitations are important.

First, the framework depends on the MRW parabolic projection family. This
makes the diagnostics interpretable, but it also means that non-MRW
multifractal processes may be only partially represented. The coordinate
\(\lambda^2_{\mathrm{proj}}\) should therefore be interpreted as an
MRW-family curvature coordinate, not as a universal intermittency parameter.

Second, the finite-sample identifiability result is conditional on the tested
sample lengths, \(q\)-grid, scale sets, and estimator family. The current
evidence does not rule out better recovery with longer samples, denser scale
support, or different estimators. It does show that short-window recovery is
not reliable under the tested setting.

Third, the reported diagnostic scores are calibrated scores in the modeling
sense, not fully validated statistical probabilities. The project does not yet
include reliability diagrams, expected calibration error, or Brier-score
calibration analysis. The scores should therefore be read as comparable
diagnostic indices rather than posterior probabilities of a data-generating
mechanism.

Fourth, the real-data evidence is exploratory. The Fama-French experiment uses
surrogate comparison to check diagnostic responsiveness, not to establish a
true data-generating mechanism. Stronger real-world claims would require
domain-specific preprocessing, additional surrogate designs, uncertainty
analysis, and comparison against alternative multifractal estimators.

Finally, the classical baseline comparison is still not exhaustive. Structure
functions, MFDFA, and compact Haar wavelet controls are included, but these
wavelet rows are not substitutes for a full production WTMM or wavelet-leader
benchmark. This should be acknowledged openly in any submission.
